% Tailored for Microsoft Research Redmond - Research Intern, Data Systems.
% Requisition: 200058258; portal: 1970393557008574.
% https://apply.careers.microsoft.com/careers/job/1970393557008574
% Sources: master/ and the long-form root main.tex; prepared 2026-10-07.
% Build from this directory:
% latexmk -pdf -jobname=Hongrui_Zhang_Microsoft_Data_Systems_2026-10-07 resume.tex

\documentclass[11pt]{article}
\usepackage[letterpaper,margin=0.65in,nohead,includefoot,footskip=18pt]{geometry}
\usepackage[T1]{fontenc}
\usepackage[utf8]{inputenc}
\usepackage{palatino}
\usepackage{microtype}
\usepackage{tabularx}
\usepackage{enumitem}
\usepackage{xcolor}
\usepackage[hidelinks]{hyperref}
\usepackage{fancyhdr}

\hypersetup{
  pdftitle={Hongrui Zhang - Microsoft Research Data Systems},
  pdfauthor={Hongrui Zhang}
}
\definecolor{sectiongray}{gray}{0.92}
\pagestyle{fancy}
\fancyhf{}
\fancyfoot[L]{\small Hongrui Zhang}
\fancyfoot[R]{\small\thepage}
\renewcommand{\headrulewidth}{0pt}
\renewcommand{\footrulewidth}{0pt}
\setlength{\parindent}{0pt}
\setlength{\parskip}{0pt}
\setlength{\tabcolsep}{0pt}
\setlist[itemize]{leftmargin=1.3em,itemsep=4pt,topsep=5pt,parsep=0pt,partopsep=0pt}
\raggedbottom
\urlstyle{same}
\emergencystretch=1em

% Section headings account for colorbox padding and stay with following text.
\newcommand{\resheading}[1]{%
  \par\addvspace{12pt}%
  \noindent\colorbox{sectiongray}{%
    \parbox{\dimexpr\linewidth-2\fboxsep\relax}{\strut\small\textbf{#1}}%
  }\par\nobreak\vspace{5pt}\nopagebreak[4]%
}

% Wrapping titles leave a separate column for dates.
\newcommand{\resumeentry}[3]{%
  \par\addvspace{8pt}%
  \begin{tabularx}{\linewidth}{@{}>{\raggedright\arraybackslash}X@{\hspace{1em}}r@{}}%
    \textbf{#1} & {\small #2}\\%
  \end{tabularx}\par\nobreak%
  \if\relax\detokenize{#3}\relax\else
    {\small #3\par}\nobreak%
  \fi\nopagebreak[4]%
}

\begin{document}

\begin{center}
  {\fontsize{24}{27}\selectfont\textbf{Hongrui (Jack) Zhang}}\\[3pt]
  {\large Computer Science Ph.D. Researcher \textbar\ Data Systems and GPU Computing}\\[5pt]
  Irvine, California\enspace\textbar\enspace 302-898-6341\enspace\textbar\enspace
  \href{mailto:jackzhang010603@gmail.com}{jackzhang010603@gmail.com}\\[2pt]
  \href{https://www.linkedin.com/in/hongruizhang63}{LinkedIn}\enspace\textbar\enspace
  \href{https://hongruizhang.com}{hongruizhang.com}
\end{center}
\vspace{-4pt}

\resheading{EDUCATION}
\resumeentry{Ph.D. in Computer Science, University of California, Riverside}
  {Sep 2024 -- Expected Jun 2028}
  {Advisor: Hung-Wei Tseng }

\resumeentry{B.S. in Computer Engineering, University of California, San Diego}
  {Sep 2019 -- Sep 2023}
  {}

\resheading{RESEARCH EXPERIENCE}

\resumeentry{GPU-Accelerated Database Systems and Matrix-Native Query Processing}
  {Jun 2026 -- Present}
  {Python, PyTorch, SQL, CUDA, cuSPARSE \textbar\ Query execution and experimental evaluation}
\begin{itemize}
  \item Developed and validated a matrix-native execution plan for TPC-H Q3, expressing joins, grouping, and aggregation using sparse matrix and tensor operations.
  \item Optimized TCUDB-style sparse joins using direct one-hot CSR construction, reducing Q4 join latency from 40.5 ms to 11.0 ms (3.7$\times$) over generic COO construction at TPC-H SF10.
  \item Implemented and SQL-validated PyTorch execution plans for all 22 TPC-H queries, profiling GPU bottlenecks on an RTX 5090 to guide hardware-aware query processing.
\end{itemize}

\resumeentry{RTSpMSpM — RT-Core-Accelerated Sparse Matrix Multiplication}
  {Sep 2023 -- Sep 2025}
  {\textbf{ISCA '25} \textbar\ First author and conference presenter \textbar\ OptiX, CUDA, C++, Python}
\begin{itemize}
  \item Designed and implemented an OptiX/CUDA library mapping sparse matrix multiplication onto NVIDIA RT Cores, achieving a 1.85$\times$ geometric-mean over cuSPARSE.
  \item Proposed RT core architectural enhancements, including fused intersection-multiplication and hardware row accumulation, achieving 3.06$\times$ simulated speedup over cuSPARSE with only 0.2{\%} estimated area overhead.
  \item Integrated the library into a graph neural network training framework using CUDA and Python, achieving more than 3$\times$ speedup on NVIDIA RTX 4090.
\end{itemize}

\resheading{PUBLICATION}
\textbf{Hongrui Zhang}, Yunan Zhang, and Hung-Wei Tseng.
``RTSpMSpM: Harnessing Ray Tracing for Efficient Sparse Matrix Computations.''
\textit{52nd International Symposium on Computer Architecture (ISCA '25)}, 2025.
\href{https://doi.org/10.1145/3695053.3731072}{\textbf{[DOI]}}
\par\nobreak\vspace{3pt}
\textit{IEEE Micro Top Picks 2026}\enspace---\enspace\textbf{Honorable Mention}\par

\resheading{TECHNICAL SKILLS}
\begin{itemize}[label={},leftmargin=0pt,itemsep=5pt]
  \item \textbf{Languages:} Python, C++, CUDA C++, SQL, C, Bash
  \item \textbf{Frameworks \& Tools:} PyTorch, CUDA, NVIDIA OptiX, cuSPARSE, NumPy, Git, Docker, Linux
  \item \textbf{GPU Profiling \& Data Systems:} Nsight Systems, Nsight Compute, PyTorch Profiler, TPC-H, sparse matrix operations, query execution
\end{itemize}

\end{document}
